\documentclass[a4paper,11pt]{article}
\usepackage[margin=1in]{geometry}
\usepackage{amsmath,amsthm,amssymb}

\newtheorem{theorem}{Theorem}
\newtheorem{corollary}[theorem]{Corollary}
\DeclareMathOperator{\ObsDiam}{ObsDiam}
\DeclareMathOperator{\osc}{osc}
\DeclareMathOperator{\TV}{TV}

\title{Observable Diameter under Metric and Mass Perturbations}
\author{}
\date{}

\begin{document}
\maketitle

\begin{abstract}
On a finite metric space, observable diameter measures the largest,
over real-valued 1-Lipschitz functions, of the smallest oscillation
remaining after a prescribed amount of mass may be discarded. We
compare this quantity when both the metric and the probability weights
change on the same finite set. A uniform metric error contributes its
size additively, while a total variation error requires a corresponding
reduction in the allowed discarded mass on the comparison side. A
two-point example attains the metric bound, and another shows why the
discarded-mass shift cannot be replaced by continuity at the same
parameter.
\end{abstract}

\section{Retained oscillation and the comparison theorem}

A real-valued observable can vary widely on a few points while remaining
nearly constant on most of the probability mass. To measure its spread
on most of a space, we first minimize its oscillation on subsets carrying
enough mass. Observable diameter then takes the largest of these
minimal spreads over all 1-Lipschitz observables.

Let $X$ be a nonempty finite set. A probability weight $\mu$ on $X$
satisfies $\mu(x)\ge0$ for $x\in X$ and
$\sum_{x\in X}\mu(x)=1$; write
$\mu(A)=\sum_{x\in A}\mu(x)$ for $A\subseteq X$.
For a function $f:X\to\mathbb R$ and a nonempty subset $A\subseteq X$,
its \emph{oscillation} on $A$ is
\[
 \osc_A f=\max_{x\in A}f(x)-\min_{x\in A}f(x).
\]
For a metric $d$ on $X$, a function $f:X\to\mathbb R$ is
1-Lipschitz with respect to $d$ if
$|f(x)-f(y)|\le d(x,y)$ for all $x,y\in X$.

For $0\le\kappa<1$, define
\begin{equation}\label{eq:definition}
 \ObsDiam(X,d,\mu;\kappa)
 =\sup_{\substack{f:X\to\mathbb R\\
                  |f(x)-f(y)|\le d(x,y)\ (x,y\in X)}}
    \ \min_{\substack{A\subseteq X\\\mu(A)\ge1-\kappa}}
      \osc_A f.
\end{equation}
Here $\kappa$ is the \emph{discarded-mass parameter}: an eligible set
must retain at least $1-\kappa$ of the mass. Eligibility includes
equality at this threshold. The eligible family contains $X$ and is
finite, so the inner minimum exists. Since $1-\kappa>0$, every
eligible set is nonempty. Also
\[
 0\le\ObsDiam(X,d,\mu;\kappa)
 \le\operatorname{diam}_d X,
 \qquad
 \operatorname{diam}_d X=\max_{x,y\in X}d(x,y),
\]
because every eligible oscillation of a 1-Lipschitz function is at most
this finite diameter. The supremum in \eqref{eq:definition} is thus a
finite real number.

For two probability weights $\mu,\nu$ on $X$, define their
\emph{total variation distance} by
\begin{equation}\label{eq:tvdef}
 \TV(\mu,\nu)=\sup_{A\subseteq X}|\mu(A)-\nu(A)|.
\end{equation}
It controls the mass change of every subset, which is the information
needed to compare the eligible sets in \eqref{eq:definition}.

\begin{theorem}\label{thm:comparison}
Let $d$ and $e$ be metrics on the nonempty finite set $X$, and let
$\mu,\nu$ be probability weights on $X$. Suppose that $\delta\ge0$ and
$0\le\tau\le\kappa<1$ satisfy
\[
 |d(x,y)-e(x,y)|\le\delta\quad(x,y\in X),
 \qquad \TV(\mu,\nu)\le\tau.
\]
Then
\begin{equation}\label{eq:comparison}
 \ObsDiam(X,d,\mu;\kappa)
 \le \ObsDiam(X,e,\nu;\kappa-\tau)+\delta.
\end{equation}
\end{theorem}

The two perturbations enter the estimate in different places. The
metric error $\delta$ increases the permitted oscillation. The mass
error $\tau$ requires the comparison set to retain at least
$1-\kappa+\tau$ under $\nu$, so that it still retains at least
$1-\kappa$ under $\mu$. This requirement explains both the shift
$\kappa-\tau$ and the condition $\tau\le\kappa$.

\section{Proof of the comparison}

\begin{proof}
We first verify a finite formula for total variation, also useful for
computing the examples. Set $h(x)=\mu(x)-\nu(x)$ for $x\in X$ and
$P=\{x\in X:h(x)>0\}$. Since $\sum_{x\in X}h(x)=0$, the number
$S=\sum_{x\in P}h(x)$ satisfies
\[
 S=-\sum_{\substack{x\in X\\h(x)<0}}h(x)
   =\frac12\sum_{x\in X}|h(x)|.
\]
For every subset $A\subseteq X$, the sum $\sum_{x\in A}h(x)$ lies
between $-S$ and $S$, and the subset $P$ attains $S$. Thus
\begin{equation}\label{eq:tvsum}
 \TV(\mu,\nu)=\frac12\sum_{x\in X}|\mu(x)-\nu(x)|.
\end{equation}
In particular, the assumed total variation bound implies
$\mu(A)\ge\nu(A)-\tau$ for every $A\subseteq X$.

Now fix any $d$-1-Lipschitz function $f:X\to\mathbb R$. Define
$g:X\to\mathbb R$ by
\begin{equation}\label{eq:envelope}
 g(x)=\min_{y\in X}\bigl(f(y)+e(x,y)\bigr)\quad(x\in X).
\end{equation}
The minimum exists because $X$ is finite and nonempty. The term $y=x$
gives $g(x)\le f(x)$. For any $x,y\in X$, the Lipschitz condition for
$f$ and the metric comparison give
\[
 f(y)+e(x,y)
 \ge f(x)-d(x,y)+e(x,y)
 \ge f(x)-\delta.
\]
Taking the minimum over $y$ yields the one-sided approximation
\begin{equation}\label{eq:onesided}
 f(x)-\delta\le g(x)\le f(x)\quad(x\in X).
\end{equation}

The function $g$ is 1-Lipschitz for $e$. To see this, fix $x,z\in X$
and choose $y_z\in X$ attaining the minimum in \eqref{eq:envelope}
for $z$. By the triangle inequality for $e$,
\[
 g(x)\le f(y_z)+e(x,y_z)
       \le f(y_z)+e(z,y_z)+e(x,z)
       =g(z)+e(x,z).
\]
Interchanging $x$ and $z$ proves
$|g(x)-g(z)|\le e(x,z)$.

For any nonempty $A\subseteq X$, \eqref{eq:onesided} implies
\[
 \max_{x\in A}f(x)\le\max_{x\in A}g(x)+\delta,
 \qquad
 \min_{x\in A}f(x)\ge\min_{x\in A}g(x).
\]
Subtracting gives
\begin{equation}\label{eq:oscillation}
 \osc_A f\le\osc_A g+\delta.
\end{equation}
Only one endpoint of the oscillation needs an additive error: the
pointwise difference $f-g$ lies between zero and $\delta$. A bound by
$2\delta$ would lose this one-sided information.

Choose a subset $A_*\subseteq X$ minimizing $\osc_A g$ subject to
$\nu(A)\ge1-(\kappa-\tau)$. Such a set exists because there are
finitely many eligible sets and $X$ is one of them. It is nonempty,
and the total variation estimate gives
\[
 \mu(A_*)\ge\nu(A_*)-\tau
 \ge1-\kappa+\tau-\tau=1-\kappa.
\]
Hence $A_*$ is also eligible for the $\mu$-minimum at $\kappa$.
Using \eqref{eq:oscillation} on this set, we obtain
\begin{align*}
 \min_{\substack{A\subseteq X\\\mu(A)\ge1-\kappa}}\osc_A f
 &\le\osc_{A_*}f\\
 &\le\osc_{A_*}g+\delta\\
 &=\min_{\substack{A\subseteq X\\\nu(A)\ge1-(\kappa-\tau)}}
       \osc_A g+\delta\\
 &\le\ObsDiam(X,e,\nu;\kappa-\tau)+\delta.
\end{align*}
The last inequality holds because $g$ is $e$-1-Lipschitz and therefore
is one of the functions in the defining supremum on the right.
This bound holds for every $d$-1-Lipschitz function $f$. Taking their
supremum only now proves \eqref{eq:comparison}.
\end{proof}

\section{Reversing the comparison}

The assumptions of Theorem~\ref{thm:comparison} are symmetric in the
pairs $(d,\mu)$ and $(e,\nu)$: the metric error is unchanged, and
\eqref{eq:tvdef} gives $\TV(\nu,\mu)=\TV(\mu,\nu)$. Applying the
theorem with the pairs interchanged gives, under the same parameter
constraints,
\begin{equation}\label{eq:reverse}
 \ObsDiam(X,e,\nu;\kappa)
 \le\ObsDiam(X,d,\mu;\kappa-\tau)+\delta.
\end{equation}
Both comparisons shift the parameter on their right side. For
$\tau>0$, they do not give an absolute-difference bound at the same
discarded-mass parameter.

\begin{corollary}\label{cor:metric}
If $\mu=\nu$ and $|d(x,y)-e(x,y)|\le\delta$ for all $x,y\in X$, then
for every $0\le\kappa<1$,
\[
 \bigl|\ObsDiam(X,d,\mu;\kappa)
       -\ObsDiam(X,e,\mu;\kappa)\bigr|\le\delta.
\]
\end{corollary}

\begin{proof}
For equal weights, \eqref{eq:tvdef} gives total variation zero.
Set $\tau=0$ in \eqref{eq:comparison} and \eqref{eq:reverse}; these
are the two inequalities defining the asserted absolute bound.
\end{proof}

\section{Sharpness and a mass threshold}

\subsection{The metric error is attained}

Let $X=\{a,b\}$, with $\mu(a)=\mu(b)=1/2$, and choose
$0\le\kappa<1/2$. Let $s>0$ and $\delta\ge0$, and define metrics by
\[
 d(a,b)=s+\delta,\qquad e(a,b)=s.
\]
The diagonal distances are zero and the distances are symmetric, so
these positive off-diagonal values define metrics. Their uniform
difference is $\delta$.

The retained-mass threshold $1-\kappa$ is greater than $1/2$.
Neither singleton is eligible, and thus every eligible set is $X$.
For $d$-1-Lipschitz $f$, its oscillation on $X$ is
$|f(a)-f(b)|\le s+\delta$. This upper bound is attained by
$f(a)=0$ and $f(b)=s+\delta$, which is $d$-1-Lipschitz. Therefore
\[
 \ObsDiam(X,d,\mu;\kappa)=s+\delta.
\]
For $e$, the bound $|f(a)-f(b)|\le s$ is attained by the values
$0$ and $s$, giving
\[
 \ObsDiam(X,e,\mu;\kappa)=s.
\]
Their difference equals $\delta$, so the additive metric constant in
Corollary~\ref{cor:metric} is attained.

\subsection{Small mass changes can cross the eligibility threshold}

Keep $X=\{a,b\}$, now with distance one between the points, and denote
this metric by $d$. Fix $\kappa=0.4$ and define weights
\[
 \mu(a)=0.6,\quad\mu(b)=0.4,
 \qquad
 \nu_\varepsilon(a)=0.6-\varepsilon,\quad
 \nu_\varepsilon(b)=0.4+\varepsilon,
 \qquad 0<\varepsilon<0.1.
\]
Formula~\eqref{eq:tvsum} gives
\[
 \TV(\mu,\nu_\varepsilon)
 =\tfrac12(\varepsilon+\varepsilon)=\varepsilon.
\]
At parameter $\kappa=0.4$, the retained-mass threshold is $0.6$.
The singleton $\{a\}$ is eligible under $\mu$ and has oscillation
zero for every function. Nonnegativity of oscillation therefore gives
\[
 \ObsDiam(X,d,\mu;0.4)=0.
\]
Under $\nu_\varepsilon$, both singletons have mass less than $0.6$:
$0.6-\varepsilon<0.6$ and $0.4+\varepsilon<0.5<0.6$.
Only $X$ is eligible. Every $d$-1-Lipschitz function has oscillation
at most one on $X$, and the values $f(a)=0$, $f(b)=1$ attain one.
Consequently,
\[
 \ObsDiam(X,d,\nu_\varepsilon;0.4)=1.
\]
The total variation tends to zero with $\varepsilon$, while this
difference of observable diameters stays equal to one. Thus arbitrarily
small changes in mass can destroy continuity at a fixed discarded-mass
parameter. The inclusion of equality in the eligibility condition is
essential to this calculation: $\{a\}$ qualifies under $\mu$ precisely
at mass $0.6$.

This jump is consistent with the shifted comparison. Set
$\tau=\varepsilon$ and $e=d$, so that $\delta=0$ and
$0\le\tau\le0.4$. The reversed bound \eqref{eq:reverse} states
\[
 1=\ObsDiam(X,d,\nu_\varepsilon;0.4)
 \le\ObsDiam(X,d,\mu;0.4-\varepsilon).
\]
On the right, the retained-mass threshold is
$1-(0.4-\varepsilon)=0.6+\varepsilon$. Both singleton masses under
$\mu$ are below this threshold, so only $X$ is eligible and the
right side equals one. The bound is therefore $1\le1$.

The examples separate the two effects in Theorem~\ref{thm:comparison}:
the metric error controls oscillation additively, whereas the mass error
changes which sets are eligible. The latter is why the theorem retains
the shift in discarded mass. All statements here concern metrics and
probability weights on one nonempty finite set.

\end{document}
